% questions.tex —— 2005 年·普通高等学校招生考试·重庆卷·理科 · 题目内容片段（无答案）
% 由 tools/generate_projects.py 依来源总账自 2005/2005_数学真题_重庆理科_试题.md 生成，勿手改（改源 MD 后重跑）

\gkpart{一、选择题（本大题共 10 小题，每小题 5 分，共 50 分）}

\begin{question}
\gkti{1}圆 $(x + 2)^2 + y^2 = 5$ 关于原点 $(0, 0)$ 对称的圆的方程为\mbox{（　　）}

\begin{gkopts}(2)
\task A．$(x - 2)^2 + y^2 = 5$
\task B．$x^2 + (y - 2)^2 = 5$
\task C．$(x + 2)^2 + (y + 2)^2 = 5$
\task D．$x^2 + (y + 2)^2 = 5$
\end{gkopts}

\end{question}
\begin{question}
\gkti{2}$\left(\dfrac{1 + \mathrm{i}}{1 - \mathrm{i}}\right)^{2005} =$\mbox{（　　）}

\begin{gkopts}(4)
\task A．$\mathrm{i}$
\task B．$-\mathrm{i}$
\task C．$2^{2005}$
\task D．$-2^{2005}$
\end{gkopts}

\end{question}
\begin{question}
\gkti{3}若函数 $f(x)$ 是定义在 $\mathbf{R}$ 上的偶函数，在 $(-\infty, 0]$ 上是减函数，且 $f(2) = 0$，则使得 $f(x) < 0$ 的 $x$ 的取值范围是\mbox{（　　）}

\begin{gkopts}(2)
\task A．$(-\infty, 2)$
\task B．$(2, +\infty)$
\task C．$(-\infty, -2) \cup (2, +\infty)$
\task D．$(-2, 2)$
\end{gkopts}

\end{question}
\begin{question}
\gkti{4}已知 $A(3, 1), B(6, 1), C(4, 3)$，$D$ 为线段 $BC$ 的中点，则向量 $\overrightarrow{AC}$ 与 $\overrightarrow{DA}$ 的夹角为\mbox{（　　）}

\begin{gkopts}(2)
\task A．$\dfrac{\pi}{2} - \arccos\dfrac{4}{5}$
\task B．$\arccos\dfrac{4}{5}$
\task C．$\arccos\left(-\dfrac{4}{5}\right)$
\task D．$-\arccos\left(-\dfrac{4}{5}\right)$
\end{gkopts}

\end{question}
\begin{question}
\gkti{5}若 $x, y$ 是正数，则 $\left(x + \dfrac{1}{2y}\right)^2 + \left(y + \dfrac{1}{2x}\right)^2$ 的最小值是\mbox{（　　）}

\begin{gkopts}(4)
\task A．$3$
\task B．$\dfrac{7}{2}$
\task C．$4$
\task D．$\dfrac{9}{2}$
\end{gkopts}

\end{question}
\begin{question}
\gkti{6}已知 $\alpha, \beta$ 均为锐角，若 $p: \sin\alpha < \sin(\alpha + \beta)$，$q: \alpha + \beta < \dfrac{\pi}{2}$，则 $p$ 是 $q$ 的\mbox{（　　）}

\begin{gkopts}(2)
\task A．充分而不必要条件
\task B．必要而不充分条件
\task C．充要条件
\task D．既不充分也不必要条件
\end{gkopts}

\end{question}
\begin{question}
\gkti{7}对于不重合的两个平面 $\alpha$ 与 $\beta$，给定下列条件：

① 存在平面 $\gamma$，使得 $\alpha, \beta$ 都垂直于 $\gamma$；

② 存在平面 $\gamma$，使得 $\alpha, \beta$ 都平行于 $\gamma$；

③ $\alpha$ 内有不共线的三点到 $\beta$ 的距离相等；

④ 存在异面直线 $l, m$，使得 $l \parallel \alpha, l \parallel \beta, m \parallel \alpha, m \parallel \beta$．

其中，可以判定 $\alpha$ 与 $\beta$ 平行的条件有\mbox{（　　）}

\begin{gkopts}(4)
\task A．1 个
\task B．2 个
\task C．3 个
\task D．4 个
\end{gkopts}

\end{question}
\begin{question}
\gkti{8}若 $\left(2x - \dfrac{1}{x}\right)^n$ 展开式中含 $\dfrac{1}{x^2}$ 项的系数与含 $\dfrac{1}{x^4}$ 项的系数之比为 $-5$，则 $n$ 等于\mbox{（　　）}

\begin{gkopts}(4)
\task A．4
\task B．6
\task C．8
\task D．10
\end{gkopts}

\end{question}
\begin{question}
\gkti{9}若动点 $(x, y)$ 在曲线 $\dfrac{x^2}{4} + \dfrac{y^2}{b^2} = 1$ ($b > 0$) 上变化，则 $x^2 + 2y$ 的最大值为\mbox{（　　）}

\begin{gkopts}(1)
\task A．$\begin{cases} \dfrac{b^2}{4} + 4, & 0 < b < 4, \\ 2b, & b \geqslant 4. \end{cases}$
\task B．$\begin{cases} \dfrac{b^2}{4} + 4, & 0 < b < 2, \\ 2b, & b \geqslant 2. \end{cases}$
\task C．$\dfrac{b^2}{4} + 4$
\task D．$2b$
\end{gkopts}

\end{question}
\begin{question}
\gkti{10}在体积为 1 的三棱锥 $A - BCD$ 侧棱 $AB, AC, AD$ 上分别取点 $E, F, G$，使 $AE:EB = AF:FC = AG:GD = 2:1$，记 $O$ 为三平面 $BCG, CDE, DBF$ 的交点，则三棱锥 $O - BCD$ 的体积等于\mbox{（　　）}

\begin{gkopts}(4)
\task A．$\dfrac{1}{9}$
\task B．$\dfrac{1}{8}$
\task C．$\dfrac{1}{7}$
\task D．$\dfrac{1}{4}$
\end{gkopts}

\end{question}
\gkpart{二、填空题（本大题共 6 小题，每小题 4 分，共 24 分）}

\begin{question}
\gkti{11}已知集合 $A = \{x \in \mathbf{R} \mid x^2 - x - 6 < 0\}$，集合 $B = \{x \in \mathbf{R} \mid |x - 2| < 2\}$，则 $A \cap B =$ \gkblank{}．

\end{question}
\begin{question}
\gkti{12}曲线 $y = x^3$ 在点 $(a, a^3)$ ($a \neq 0$) 处的切线与 $x$ 轴、直线 $x = a$ 所围成的三角形的面积为 $\dfrac{1}{6}$，则 $a =$ \gkblank{}．

\end{question}
\begin{question}
\gkti{13}已知 $\alpha, \beta$ 均为锐角，且 $\cos(\alpha + \beta) = \sin(\alpha - \beta)$，则 $\tan\alpha =$ \gkblank{}．

\end{question}
\begin{question}
\gkti{14}$\lim_{n \to \infty} \dfrac{2^{3n} - 3^{2n+1}}{2^{3n} + 3^{2n}} =$ \gkblank{}．

\end{question}
\begin{question}
\gkti{15}某轻轨列车有 4 节车厢，现有 6 位乘客准备乘坐，设每一位乘客进入每节车厢是等可能的，则这 6 位乘客进入各节车厢的人数恰好为 0, 1, 2, 3 的概率为 \gkblank{}．

\end{question}
\begin{question}
\gkti{16}连接抛物线上任意四点组成的四边形可能是 \gkblank{}．（填写所有正确选项的序号）

① 菱形；② 有 3 条边相等的四边形；③ 梯形；④ 平行四边形；⑤ 有一组对角相等的四边形．

\end{question}
\gkpart{三、解答题（本大题共 6 小题，共 74 分）}

\begin{question}
\gkti{17}（本小题满分 12 分）

若函数 $f(x) = \dfrac{1 + \cos 2x}{4\sin\left(\dfrac{\pi}{2} + x\right)} - a \sin\dfrac{x}{2} \cos\left(\pi - \dfrac{x}{2}\right)$ 的最大值为 2，试确定常数 $a$ 的值．

\end{question}
\begin{question}
\gkti{18}（本小题满分 12 分）

在一次购物抽奖活动中，假设某 10 张券中有一等奖券 1 张，可获价值 50 元的奖品；有二等奖券 3 张，每张可获价值 10 元的奖品；其余 6 张没有奖，某顾客从此 10 张券中任抽 2 张，求：

(1) 该顾客中奖的概率；

(2) 该顾客获得的奖品总价值 $\xi$（元）的概率分布列和期望 $\mathrm{E}\xi$．

\end{question}
\begin{question}
\gkti{19}（本小题满分 12 分）

已知 $a \in \mathbf{R}$，讨论函数 $f(x) = \mathrm{e}^x(x^2 + ax + a + 1)$ 的极值点的个数．

\Needspace{0.4\textheight}
\end{question}
\begin{question}
\gkti{20}（本小题满分 12 分）

如图，在三棱柱 $ABC - A_1B_1C_1$ 中，$AB \perp$ 侧面 $BB_1C_1C$，$E$ 为棱 $CC_1$ 上异于 $C, C_1$ 的一点，$EA \perp EB_1$．已知 $AB = \sqrt{2}, BB_1 = 2, BC = 1, \angle BCC_1 = \dfrac{\pi}{3}$，求：

(1) 异面直线 $AB$ 与 $EB_1$ 的距离；

(2) 二面角 $A - EB_1 - A_1$ 的平面角的正切值．
\nopagebreak[4]
\begin{center}\begin{tikzpicture}[>=stealth, scale=1.3, line join=round, line cap=round, thick]
  % Coordinates
  \coordinate (A) at (0, 2.5);
  \coordinate (A1) at (3.5, 2.5);
  \coordinate (B) at (0, 0.6);
  \coordinate (B1) at (3.5, 0.6);
  \coordinate (C) at (-1.1, 0.2);
  \coordinate (C1) at (2.4, 0.2);
  \coordinate (E) at ($(C)!0.5!(C1)$);

  % Dashed lines (hidden / interior)
  \draw[dashed] (A) -- (B);
  \draw[dashed] (B) -- (C);
  \draw[dashed] (B) -- (B1);
  \draw[dashed] (B) -- (E);
  \draw[dashed] (E) -- (B1);
  \draw[dashed] (A) -- (B1);

  % Solid lines (outer / visible)
  \draw (A) -- (A1) -- (B1);
  \draw (A) -- (C) -- (C1) -- (B1);
  \draw (A1) -- (C1);
  \draw (A) -- (E);
  \draw (A1) -- (E);

  % Labels
  \node[above left] at (A) {$A$};
  \node[above right] at (A1) {$A_1$};
  \node[above left] at (B) {$B$};
  \node[right] at (B1) {$B_1$};
  \node[below left] at (C) {$C$};
  \node[below] at (E) {$E$};
  \node[below right] at (C1) {$C_1$};
\end{tikzpicture}
\end{center}
\nopagebreak[4]
\end{question}
\begin{question}
\gkti{21}（本小题满分 12 分）

已知椭圆 $C_1$ 的方程为 $\dfrac{x^2}{4} + y^2 = 1$，双曲线 $C_2$ 的左、右焦点分别为 $C_1$ 的左、右顶点，而 $C_2$ 的左、右顶点分别是 $C_1$ 的左、右焦点．

(1) 求双曲线 $C_2$ 的方程；

(2) 若直线 $l: y = kx + \sqrt{2}$ 与椭圆 $C_1$ 及双曲线 $C_2$ 都恒有两个不同的交点，且 $l$ 与 $C_2$ 的两个交点 $A$ 和 $B$ 满足 $\overrightarrow{OA} \cdot \overrightarrow{OB} < 6$（其中 $O$ 为原点），求 $k$ 的取值范围．

\end{question}
\begin{question}
\gkti{22}（本小题满分 14 分）

数列 $\{a_n\}$ 满足 $a_1 = 1$ 且 $a_{n+1} = \left(1 + \dfrac{1}{n^2 + n}\right) a_n + \dfrac{1}{2^n}$ ($n \geqslant 1$)．

(1) 用数学归纳法证明：$a_n \geqslant 2$ ($n \geqslant 2$)；

(2) 已知不等式 $\ln(1 + x) < x$ 对 $x > 0$ 成立，证明：$a_n < \mathrm{e}^2$ ($n \geqslant 1$)，其中无理数 $\mathrm{e} = 2.71828\cdots$．

\end{question}